\section{讨论}
\label{sec:discussion}

腹膜侧能站住的是样本中位数，不是空间邻接，也不是下游功能。C3上皮或成纤维乘巨噬C3AR1，在3份对25份的标签置换里q为0.027。这个q只覆盖事后点名的7条，而且两条C3共用同一个受体。SPP1巨噬乘CD44巨噬的q是0.49，不能和C3写成同一条已经增强的通讯。CCL2若从上皮接到T细胞的CCR2，q也是0.027，但乘积只有0.021；从成纤维接到巨噬CCR2则没有升高。TGFB1的乘积来自T细胞。

这些样本来自已经发表的胃癌单细胞队列。这里的置换只说明在这份公开数据里，3份腹膜肿瘤的C3乘积位于原发样本的高侧。巨噬STAT3转录均值也偏高，SOCS3没有一起升高，不能写成STAT3活性已经上来。GSE308231的3对3没有复现这个乘积：成纤维C3乘巨噬C3AR1的腹膜中位数更低，7条q都高于0.35。它不增加已经通过检验的样本量，也不代替蛋白邻近或受体阻断。

原发灶的保留规则收成0对。APP--CD74等3对达到5/9例，界面中位数也明显高于事先指定的5对，但留一后过线数下降，不能写成在9例上都站得住的边。事先指定的TGFB1--TGFBR2最多3例过线，中位数仍为0，和“至少5例且留一稳定”之间差了留一和切片数两道。

接触比例低于空分布，和界面信号里大量0中位数是同一几何条件的两面。150\,$\mu$m已经盖住Visium相邻点位，前25\%又使成纤维高分区足够密，随机标签的接触比例被推到0.97附近。观察值低于这条线，说明在这个定义下看不到超出随机的邻域富集。胶原、纤连蛋白相关边的绝对信号高，是因为配体和受体在很多点位上都有表达；坐标打乱后信号仍高，所以绝对值和空间超额是两件分开的事。

腹膜样本均值里，C3、SPP1、CCL2在3份腹膜肿瘤上都高于原发中位数，配对患者上也同向。空间检验里，C3--C3AR1和SPP1--CD44只在1/9例原发灶上超过坐标打乱，CCL2--CCR2为0/9。表达方向和界面超额因此没有落在同一批边上。TGFB1的两组中位数几乎相同，空间上TGFB1--TGFBR2也只有3例过线。3份腹膜样本不能把这种方向写成腹膜特异的阻断顺序。

把中位数换成比例加权、把坐标打乱换成配体图平移，绝对读数不再全是0，超额仍然站在空分布下面。这说明第一套的阴性结果不单是中位数被零压住。半径从100改到300\,$\mu$m，方向不变。GC4在旧空分布下的那一个C3--C3AR1过线，没有在平移空分布下变成9例的超额。

单细胞分类把三条腹膜方向拆开了。C3在成纤维和上皮里都高于原发中位数，SPP1主要在巨噬，CCL2的腹膜样本均值并不是成纤维升高。3份对26份仍只是方向。
